
%============================= HEADER =================================
\chapter{The Interstellar Medium}\label{ch:i}

For this part of this chapter, I'll be mainly following \cite{2011piim.book.....D}, Chapters 1-2, but I won't fail to report other used references were I to do so.
\par The interstellar medium, or ISM, is, arguably, also the most important component of galaxies, for it is the ISM that is responsible for forming the stars that are the dominant sources of energy. While it now appears that the mass of most galaxies is primarily in the form of dark matter particles that are collisionless, or nearly so, it is the baryons (accounting for perhaps $\sim 10\%$ of the total mass) that determine the visible appearance of galaxies, and that are responsible for nearly all of the energy emitted by galaxies, derived from nuclear fusion in stars and the release of gravitational energy in accretion disks around black holes.
\par The ISM has mainly four constituents:
\begin{enumerate}
    \item Gas, which covers a whooping $99\%$ of the total mass of the ISM. Note that when we refer to gas, we're imagining it as a macroscopic fluid/body;
    \item Dust \marginpar{With the term refractory elements we're referring to solid particles that evaporate/sublimate at temperatures $T \gg 1000\unit{K}$. For the sake of simplicity, they're usually assumed to be spherical.}, which has a significantly lower density than gas $\rho_{d} \sim 1\% \rho_{g}$. Dust is made of a hefty amount of \emph{refractory elements}, mainly coming in the form of grains of different morphology and diameters in the range $D \sim 0.1 - 10 \,\mu\text{m}$. Dust is mostly present in cold environments, but incidentally also in regions where $T \sim 10^{5}-10^{6}$;
    \item Cosmic Rays are charged, highly energetic particles produced in extreme environments, like supernovae and the stellar winds of O/B-class stars. The vast majority of CR are protons ($\sim 90\%$), followed by $\alpha$ particles and electrons. Both have velocities much greater than those typical of the ISM: The most energetic ever detected had an energy of $\sim 10^{21}\unit{eV}$. As such, they can easily penetrate even the most dense regions of the ISM.
    \item Photons from many sources, including the cosmic microwave background (CMB) in the millimeter/radio band, stellar photospheres (i.e., starlight) in the optical band, radiation emitted by interstellar ions, atoms, and molecules in the infrared band; thermal emission from interstellar grains that have been heated by starlight still in the infrared. Then moving to the higher energy ($\gamma$ and X) we also observe photons from free-free emission (“Bremsstrahlung”) from interstellar plasma, synchrotron radiation from relativistic electrons, and gamma rays emitted in nuclear transitions and $\pi_{0}$ decays.
\end{enumerate}
The mass of the gas component alone makes about $10\%$ of the total mass of a galaxy, but it still is lower than the mass "stored" through stars. That said, the ISM is far more spatially extended than stars ($\sim 25-27\unit{kpc}$). Typical conversion rates gas$\longleftrightarrow$stars are 
\begin{equation*}
    \der{M}{t}(\text{gas}\to\text{stars}) \sim 1-5M_{\odot}\unit{yr}^{-1} \quad \der{M}{t}(\text{stars}\to\text{gas}) \sim 1M_{\odot}\unit{yr}^{-1}
\end{equation*}

The baryons in the interstellar medium of the Milky Way are found with a wide range of temperatures and densities; because the interstellar medium is dynamic, all densities and temperatures within these ranges can be found somewhere in the Milky Way. However, it is observed that most of the baryons have temperatures falling close to various characteristic states, or “phases.” For purposes of discus- sion, it is convenient to name these phases. Here we identify seven distinct phases\footnote{Although how well-separated these phases are is a completely different question.} that, between them, account for most of the mass and most of the volume of the interstellar medium. The seven phases are defined based on the temperature, density, and ionization status of Hydrogen atoms. These phases consist of the following:
\begin{enumerate}
    \item \emph{Hot ionized medium} (HIM) or \emph{Coronal Gas}:  Gas that has been shock-heated to temperatures $T>10^{6}\unit{K}$ by blastwaves racing outward from supernova explosions. The gas is collisionally ionized, with ions such as OVI (i.e. $O^{5+}$) present. Most of the coronal gas has low density ($n<0.01\unit{cm}^{-3}$), filling an appreciable fraction ($10-50\%$) of the volume of the galactic disk. The coronal gas regions may have characteristic dimensions of $\sim 20 \unit{pc}$, and may be connected to other coronal gas volumes. The coronal gas cools on $\sim\unit{Myr}$ time scales. Much of the volume above and below the disk is thought to be pervaded by coronal gas. The main heating mechanisms are (in individual galaxies) supernovae blast waves, (in galaxy clusters) release of gravitational energy, while they cool through X/$\gamma$-rays emission. The main tracers for this region are X/$\gamma$/UV-rays emission and radio emission from synchrotron radiation;
    \item \emph{Warm Ionized} H$_{\text{II}}$ \emph{Medium}: \marginpar{The Sun is located in the (warm) Local Bubble, a cavity with diameter of $\sim300\unit{lyr}$, which could have been created by the esplosion of a nearby Supernova (Geminga),$300,000\unit{yr}$ ago} Gas where the hydrogen has been photoionized by ultraviolet photons from hot stars. Most of this photoionized gas is maintained by radiation from recently formed hot massive O-type stars - the photoionized gas may be dense material from a nearby cloud (in which case the ionized gas is called an H$_{\text{II}}$ region) or lower density “intercloud” medium (referred to as diffuse H$_{\text{II}}$). The density of the WIM depends on the particular environment (but roughly $n \sim 1\unit{cm}^{-3}$), but has a temperature of about $T\sim 10^{4}-10^{5}\unit{K}$ and fills about $10-15\%$ of the volume of the galactic disk. The main heating mechanism is UV radiation from O-B stars, while they cool through high-excitation transitions of atoms. The main tracers for this region are radio lines and continuum (recombination lines, free-free), optical/UV lines;
    \item \emph{Warm neutral medium} (WNM): Predominantly atomic gas heated to temperatures $T \approx 10^{3.7}\unit{K}$. In the local interstellar medium, this gas is found at densities $n \approx 0.6\unit{cm}^{-3}$. It fills a significant fraction of the volume of the disk, perhaps 40\%. The main heating mechanism is photoelectric effect on grains by photons, while cooling occurs mainly through the $21\unit{cm}$ line of atomic Hydrogen or through the $158\,\mu\text{m}$ line of ionized Carbon, which are also the main tracers;
    \item \emph{Cold neutral medium} (CNM): Predominantly atomic gas at temperatures $T \approx 10^{2}\unit{K}$, with densities $n \approx 30\unit{cm}^{-3}$ filling $\sim1\%$ of the volume of the local interstellar medium. The main heating mechanisms are the interaction with cosmic rays and photoelectric effect on grains by photons, while cooling occurs mainly through the $21\unit{cm}$ line of atomic Hydrogen, which is also the main tracer;
    \item \emph{Molecular gas}: Two different phases can be distinguished here: the \emph{diffuse molecular gas} and the \emph{dense molecular gas}.
    \begin{itemize}
        \item \emph{Diffuse molecular gas}: similar to the cool H$_{\text{I}}$ clouds, but with sufficiently large (column) densities so that the self-shielding of H$_{2}$ allows H$_{2}$ molecules to be abundant in the cloud interior. It has a temperature of $T\sim 50\unit{K}$ and a density $n \sim 100\unit{cm}^{-3}$ and fills about $0.1\%$ of the total galactic volume. The main heating mechanism is the interaction with cosmic rays, while cooling occurs mainly through the lines of atomic metals. The main tracer is molecular Hydrogen (but in absorption alone);
        \item \emph{Dense molecular gas}: Gravitationally bound clouds achieving $n > 100\unit{cm}^{-3}$. These clouds are often “dark” - with visual extinction $A_{V} \gtrsim 3 \unit{mag}$ through their central regions. In these dark clouds, the dust grains are often coated with “mantles” composed of H$_{2}$O and other molecular ices. It is within these regions that star formation takes place. They fill only the $0.001\%$ of the total galactic volume, but make about $1/4$ of the total mass. The main heating mechanism is the interaction with cosmic rays, while cooling occurs mainly through the roto-vibrational molecular lines of molecules, which are also the main tracers.
    \end{itemize}
\end{enumerate}
\begin{figure}[h!]
    \centering 
    \includegraphics[width=0.6\linewidth]{img/hydrogendistrib.png}
    \caption{Distribution of different ionization states of Hydrogen as a function of the galactocentric radius. Credits: Unknown.}
    \label{fgr:hydrogendistrib}
\end{figure}

\section{Collisional Processes}

Collisions are fundamental to the physics of the interstellar medium (ISM): They allow the gas to (usually) be treated as a fluid; they determine the thermal and electrical conductivity and diffusion coefficients; they produce most of the excitations of ions, atoms, and molecules that result in emission of photons from the ISM; and they are responsible for recombination of electrons and ions, and for chemical reactions.

\subsection{Collisional Rate Coefficients}

The rate per unit volume of a general two-body collisional process 
\begin{equation*}
    \text{A} + \text{B} \to \text{products}
\end{equation*}
is written as
\begin{equation}
    r_{AB} = n_{A}n_{B}\langle \sigma v \rangle_{AB}
\end{equation}
where the two-body collisional rate coefficient for the general collisional process is 
\begin{equation}
    \langle \sigma v \rangle_{AB} = \int_{0}^{\infty} \sigma_{AB}(v)vf_{f}\diff{v}
    \label{eq:collisionalratecoeff}
\end{equation}
where $\sigma_{AB}$ is the velocity velocity-dependent reaction cross section for the reaction and we're allowing the individual particles to have a non-trivial velocity distribution $f_{v}\diff{v}$. The collision rate coefficient thus includes two contributions: An “internal” part that takes into account both the phyisico/chemical properties of the colliders, and the quantum probabilities that a given interaction will produce a given product, and an “external” part that takes into account the kinetical status (i.e. $v$) of colliders.
\par If we assume thermal equilibrium at temperature $T$, the distribution function for the relative speed in encounters between particles A and B is given by a Maxwellian velocity distribution
\begin{equation}
    f_{v}\diff{v} = 4\pi \left(\frac{\mu}{2\pi kT}\right)^{3/2}\exp\left(-\frac{\mu v^{2}}{2kT}\right)v^{2}\diff{v}
    \label{eq:maxwelliandistrib}
\end{equation}
where $\mu$ is the reduced mass of the collision partners. Note that if the velocity distribution of the individual collision partners is a Maxwellian distribution, the (relative) velocity distribution is still a Maxwellian and has the form (\ref{eq:maxwelliandistrib}).
\par Let us give a brief overview of the collisional rate coefficient for different types of interaction.

\subsubsection{Interaction between neutral species}

The main forces at play between atoms, grains, and neutral molecules are long-range forces (attractive forces between istantaneous dipoles) and close-range forces (repulsive Coulombian forces of electron). This configuration is easily approximated by the interaction between rigid spheres, allowing us to find an expression with the likes of 
\begin{equation}
    \langle \sigma v \rangle_{nn} = 1.81\power{10}{-10}\left(\frac{T}{100\unit{K}}\right)^{1/2}\left(\frac{m_{\text{H}}}{\mu}\right)^{1/2}\left(\frac{r_{A}+r_{B}}{2\unit{Å}}\right)^{2}\unit{cm}^{-3}\unit{s}^{-1}
\end{equation} 

\subsubsection{Interaction between neutral and charged species}

This is the case of an atom or neutral molecule interacting with an ion or electron. The relevant forces at play here are the attractive dipolar forces, which are non-negligible even at long distances. Unfortunately, this means that the rigid sphere approximation is no longer valid.
\par We should then consider the force exerted by the ion on the induced dipole
\begin{equation}
    F = - \frac{2\alpha Z^{2}e^{2}}{r^{5}}
\end{equation}
to which is associated an interaction potential 
\begin{equation}
    U(r) = -\frac{1}{2}\frac{\alpha Z^{2}e^{2}}{r^{4}}
\end{equation}
A proper discussion should then consider the impact parameter for the process and deduce from there the collisional rate coefficient, which we'll here present without proof
\begin{equation}
    \langle \sigma v \rangle_{nc} = 8.98\power{10}{-10}Z\left(\frac{\alpha}{a_{0}^{3}}\right)^{1/2}\left(\frac{m_{\text{H}}}{\mu}\right)^{1/2}\unit{cm}^{-3}\unit{s}^{-1}
\end{equation}
Notably, it is independent of the temperature! So, in a cold gas, neutral-charged interactions have a higher collisional rate coefficient than the neutral-neutral scenario.

\section{Statistical Mechanics}

The interstellar medium (ISM) and intergalactic medium (IGM) are generally far from thermodynamic equilibrium. Nevertheless, the methods of statistical mechanics and thermodynamics provide powerful tools for understanding the non-equilibrium conditions that prevail, and for relating forward and reverse rates for the processes, such as ionization and recombination, that shape the medium.
\par But before getting any deeper than this, we should ask whether it is sensible of us to assume the ISM like a fluid. The general definition of fluid--which doesn't trouble the poor Knudsen and his number--is that of a macroscopic body with a bulk velocity. To achieve such a specific condition, the system we're trying to probe must have a negligible stochastic dispersion with respect to the bulk velocity on a certain scale. We'll look into those scales that satisfy $\lambda \gg l = 1/n\sigma$, where $l$ is the mean free path of the particles gathered in the wannabe fluid. As a rule of thumb, you might want to use the following expression 
\begin{equation}
    l = \frac{1}{\sigma} = 55\left(\frac{r}{1\unit{Å}}\right)^{-2}\left(\frac{n}{1\unit{cm}^{-3}}\right)^{-1}\unit{AU}
\end{equation}

\subsection{Partition Function and Levels Population}

Consider some physical system (e.g., atoms in a box of volume $V$) that is able to exchange energy with a “heat reservoir” at temperature $T$. The theory of statistical mechanics defines the partition function $\mathcal{Z}$ to be \marginpar{Note that we're implicitly assuming the particles to be identical and the system to be isolated.}
\begin{equation}
    \mathcal{Z} = \sum_{i}g_{i}e^{-E_{i}/kT}
\end{equation}
We know that the total number of particles $N$ (which is fixed) must be equal to the sum of all the particles in the different energy states
\begin{equation}
    N = \sum_{i}N_{i}
\end{equation}
It is fairly easy, although tedius, to show that the number of particles in the $i$-th energy state is given by 
\begin{equation}
    N_{i} = \frac{N}{\mathcal{Z}(T)}g_{i}e^{-E_{i}/kT}
\end{equation}
This configuration is crucial for the establishing of Local Thermodynamic Equilibrium: Matter obeys to the Boltzmann distribution for energy and Maxwell distribution for velocity, but LTE means that matter is not in \emph{global} equilibrium with radiation interacting with it.
\par What if now we relax the condition of the particles being identical? The process grows immediately more tedious and verbose, but the procedure to follow is essentially the same.
\par If we get an isolated system with total energy $E$, total mass $\mathcal{M}$, with $M$ reactant species, $N$ product species: the most likely macroscopic state is that with maximum number of compatible microscopic states, that is the one that maximizes entropy. This time we need to find the ways in which the particles of each species can be combined.
\par Skipping all the calculations we find, instead of the number of particles, the volume densities that maximize the entropy
\begin{equation}
    \boxed{
    \frac{\prod_{i=1}^{N}n_{P_{i}}^{b_{i}}}{\prod_{j=1}^{M}n_{R_{j}}^{b_{j}}} = \left[\frac{(2\pi kT)^{3/2}}{h^{3}}\right]^{\sum a_{i} - \sum b_{j}}\left[\frac{\prod_{i = 1}^{N}M_{P_{i}^{b_{i}}}}{\prod_{j = 1}^{M}M_{R_{j}^{b_{j}}}}\right]^{3/2}\frac{\prod_{i=1}^{N}z_{\text{int}}(P_{i}, T)^{b_{i}}}{\prod_{j=1}^{M}z_{\text{int}}(R_{j}, T)^{a_{j}}}
    }
    \label{eq:lawmassaction}
\end{equation}
which is the \emph{law of mass action}. At equilibrium, the number of particles per unit volume ($n=N/V$) of each species involved
is related only to temperature and to the partition functions of the individual species. The law is valid for every reaction or collision that happens in LTE.
\par Assume a reaction 
\begin{equation*}
    R_{1}+R_{2}+...+R_{M} \leftrightharpoons P_{1}+P_{2}+...+P_{N}
\end{equation*}
so that the ratio between the collision rate coefficients is 
\begin{equation*}
    \frac{\langle \sigma v \rangle_{\leftarrow}}{\langle \sigma v \rangle_{\rightarrow}} = \frac{\prod_{i=1}^{M}n_{R_{i}}}{\prod_{j=1}^{N}n_{P_{j}}}
\end{equation*}
so that plugging in the law of mass action \marginpar{Perhaps I'm tripping (fairly possible), but the law of detailed balance seems to be just the reciprocal of (\ref{eq:lawmassaction}).}
\begin{equation}
    \boxed{
    \left(\frac{\langle \sigma v \rangle_{\leftarrow}}{\langle \sigma v \rangle_{\rightarrow}}\right)^{-1} = \left[\frac{(2\pi kT)^{3/2}}{h^{3}}\right]^{N - M}\left[\frac{\prod_{i = 1}^{N}M_{P_{i}^{b_{i}}}}{\prod_{j = 1}^{M}M_{R_{j}^{b_{j}}}}\right]^{3/2}\frac{\prod_{i=1}^{N}z_{\text{int}}(P_{i}, T)^{b_{i}}}{\prod_{j=1}^{M}z_{\text{int}}(R_{j}, T)^{a_{j}}}
    }
    \label{eq:detailedbalance}
\end{equation}
results in the law of detailed balance.
\par This is fairly useful when applied to a realistic scenario, such as that of collisional (de)excitation
\begin{equation*}
    X(l) + Y \leftrightharpoons X(u) + Y
\end{equation*}
which has 
\begin{equation}
    \left(\frac{\langle \sigma v \rangle_{ul}}{\langle \sigma v \rangle_{lu}}\right) = \frac{z_{\text{int}}(X(l), T)}{z_{\text{int}}(X(u), T)} \implies \left(\frac{\langle \sigma v \rangle_{ul}}{\langle \sigma v \rangle_{lu}}\right) = \frac{g_{l}}{g_{u}}\exp\left(E_{ul}/kT\right)
\end{equation}
The rate coefficients for a process of (de-)excitation of a given species are regulated by: (i) Temperature, (ii) degeneracy of involved levels, (iii) their energy difference.
\par Another example might be that of a generic chemical reaction
\begin{equation*}
    X + Y \leftrightharpoons Z + W
\end{equation*}
which has 
\begin{equation}
    \frac{\langle \sigma v \rangle_{\leftarrow}}{\langle \sigma v \rangle_{\rightarrow}} = \frac{z_{\text{int}}(X, T)z_{\text{int}}(Y, T)}{z_{\text{int}}(Z, T)z_{\text{int}}(W, T)} \propto \frac{\exp\left(-(E_{X}+E_{Y})/kT\right)}{\exp\left(-(E_{Z}+E_{W})/kT\right)}
\end{equation}

\section{Energy Levels of Molecules}

To properly discuss the formation and "position" of the energy levels of a given molecule, we should consider, in principle, quite a lot of contributions: The kinetic term, the attractive potentials of each atom, the electronic repulsion, and the electrostatic repulsion between the atoms themselves. The resulting Hamiltonian is rather complicated to treat as it is.
\par We then introduce the \emph{Born-Oppenheimer approximation}, which assumes that Nuclei are essentially fixed and the electrons are free to move: The result is that the motion of the nuclei does not induce electronic transitions due to the different mass of electrons and nuclei. Under this approximation, you can demonstrate \marginpar{$E_{el}$ is the energy of electronic state, $E_{vib}$ energy due to radial vibration of nuclei with respect to the equilibrium position, $E_{rot}$ energy of rotation of nuclei about the
centre of mass.}
\begin{equation*}
    E_{tot} = E_{el} + E_{vib} + E_{rot}
\end{equation*}
As a further approximation, we'll assume that centrifugal distortion is negligible and only symmetric stretching is allowed.
\begin{figure}[h!]
    \centering
    \includegraphics[width=0.5\linewidth]{img/molecularmotion.png}
    \label{The possible roto-vibrational motions of a molecule.}
\end{figure}
\par Like in atoms, also in molecules the electronic transitions emit in the optical/UV bands, while the roto-vibrational transitions fall in the radio/mm/infrared bands. Note, however, that at the temperature of the molecular gas ($T<100\unit{K} ~10^{-2}\unit{eV}$) only the rotational levels in the ground vibrational and electronic states are populated!
\par The relevant parameters we'll need in our discussion are
\begin{itemize}
    \item The classical energy $E$ of a rigid rotator 
    \begin{equation*}
        E = \frac{J_{a}^{2}}{2I_{a}}+\frac{J_{b}^{2}}{2I_{b}}+\frac{J_{c}^{2}}{2I_{c}} \quad J^{2} = J_{a}^{2}+J_{b}^{2}+J_{c}^{2}
    \end{equation*}
    where $a,\,b$ and $c$ are the principal axis of the rotator;
    \item The permanent dipole moment $\mu$. If the dipole moment is non-zero, electric dipole transitions are allowed. On the other hand, if the dipole moment is vanishing, only electric quadrupole transitions (and onwards) are allowed. These, however, are strongly disfavored.
\end{itemize}
In the following, we'll consider these configurations 
\begin{itemize}
    \item Diatomic, linear: $I_{a} = I_{b} \neq I_{c} = 0$;
    \item Symmetric-top: $I_{a} = I_{b} \neq I_{c} \neq 0$;
    \item Asymmetric-top: $I_{a}\neq I_{b} \neq I_{c}$;
    \item Spherical: $I_{a} = I_{b} = I_{c}$
\end{itemize}

\subsubsection{Linear Rotors}

The classical energy is easily expressed in terms of the total (square) angular momentum $E = J^{2}/2I$, which provides us a simple expression for the quantum version. Hence 
\begin{equation*}
    E_{j} = B_{\text{rot}}hj(j+1) \quad B_{\text{rot}} = \frac{h}{8\pi^{2}I}
    \label{eq:linearrotor}
\end{equation*}
and selection rules force transitions with $\Delta j = \pm 1$. Note that from (\ref{eq:linearrotor}) we deduce that $E_{j}\propto j^{2}$ and $h\nu = E_{j+1}-E_{j} \propto j$. Also note that $E_{j} \propto B_{\text{rot}} \propto I^{-1}$, hence species with larger $B_{\text{rot}}$ (at fixed $j$) have larger $E_{j+1},\,E_{j}$ and $\nu$.
\par This has a rather troublesome implication: Molecular Hydrogen is only seen through absorption lines, and rarely from emission lines. This is for two reasons:
\begin{enumerate}
    \item Molecular Hydrogen is \emph{omonuclear}, so no dipole transitions but only electric quadrupole transitions $\Delta j = \pm 2$;
    \item Molecular Hydrogen has \emph{minimum moment of inertia}, so a huge separation between levels, and high $T$ is needed to populate the excited states. Roto-vibrational transitions are visible only in very hot gas.
\end{enumerate}
The lowest energy quadrupole transition for molecular Hydrogen has $\Delta E = 510 \unit{K}$.

\subsubsection{Symmetric-top Rotors}

For the sake of simplicity, let us define $I_{a} = I_{b} = I_{\perp}$, and $I_{c} = I_{\parallel}$. The classical energy is thus written as 
\begin{equation*}
    E = \frac{J^{2}}{2I_{\perp}} + J_{c}^{2}\left(\frac{1}{2I_{\parallel}} - \frac{1}{2I_{\perp}}\right)
\end{equation*}
To the two quantum operators $J^{2}$ and $J_{c}$ are associated the eigenvalues $j(j+1)\hbar^{2}$ and $m\hbar$, with $m \leq |j|$. We can define two rotational constants 
\begin{equation*}
    A = \frac{h}{8\pi^{2}I_{\parallel}} \quad B = \frac{h}{8\pi^{2}I_{\perp}}
\end{equation*}
which we can use to write the rotational energy
\begin{equation}
    E_{j,m} = Bhj(j+1) + (A-B)hm^{2} \quad h\nu = 2hBj
    \label{eq:symmetrictoprotors}
\end{equation}
and selection rules force transitions with $\Delta j = \pm 1$ and $\Delta m = 0$.

\subsubsection{Asymmetric-top Rotors}

This case is somewhat complicated and simple at the same time, because a general analytical expression that holds for all asymmetric-top rotators for $E_{j}$ and $\nu_{ij}$ does not exist. The rotational energy will depend, not surprisingly, on two more quantum numbers, and we have to define three rotational constants.

\begin{tcolorbox}[enhanced,attach boxed title to top center={yshift=-3mm,yshifttext=-1mm},
  colback=blue!5!white,colframe=blue!75!black,colbacktitle=blue!80!black,
  title=Example: H$_{2}$O,
  boxed title style={size=small,colframe=blue!50!black}]
  The molecule of water has the two atoms of Hydrogen form an angle of 104.45°. It has a blatant axis of symmetry passing from the atom of Oxygen, thus meaning the molecule has three different momenta of inertia.
  \par It is crucial to distinguish two possible configurations: \emph{ortho} (the two Hydrogen atoms have parallel spins) and para water (the two Hydrogen atoms have anti-parallel spins). The complete wave function for the molecule of water (which is antisymmetric) is the product of four contributions $\psi = \psi_{el}\psi_{vib}\psi_{rot}\psi_{s}$, where the former two are symmetric, while the latter two must always have opposing parity.
  \begin{itemize}
    \item Ortho: $m_{-1}, m_{+1} = $ even, odd (eo) or odd, even;
    \item Para: $m_{-1}, m_{+1} = $ even, even (ee) or odd, odd;
  \end{itemize}
  but since selection rules force $\Delta j = 0,\pm 1$ and $\Delta m = \pm 1, \pm 3$, only transitions eo $\to$ oe or ee$\to$oo (and their inverse) are allowed, while transitions between ortho and para states are strongly forbidden.
\end{tcolorbox}

\subsubsection{Spherical Rotors}

The expression this time is fairly simple
\begin{equation}
    E = \frac{J^{2}}{2I} \implies E_{j} = \frac{j(j+1)\hbar^{2}}{2I} \quad h\nu = \frac{(j+1)\hbar^{2}}{I}
\end{equation}
Note that non-polar molecules have no rotational dipolar lines unless they're distorted. Distortion, however, is only relevant for high values of $j$.

\subsection{Lines Hyperfine Structure}

Some nuclei have non-zero nuclear spin, which interacts (either via magnetic or electric quadrupole momentum) with electrons and/or the other nuclear spins. The energetic splitting due to this interaction gives the \emph{line hyperfine structure}.
Only some nuclei have non-zero nuclear spin, indicated usually with quantum number I, which depends on the particular combination of number of protons and neutrons in the nucleus.
For example, the molecule C$\element{O}{17}{}$ has the $\element{O}{17}{}$ nucleus with nuclear spin $I = 5/2$ and presents the splitting shown in Fig.\ref{fgr:c17o}.
\begin{figure}[h!]
    \centering 
    \includegraphics[width=0.6\linewidth]{img/c17o.png}
    \caption{Hyperfine splitting induced by nuclear spin on the C$\element{O}{17}{}$ molecule.}
    \label{fgr:c17o}
\end{figure}

\subsection{Inversion Transitions of NH$_{3}$}

To address what's commonly known as inversion transition, as far as the molecule of Ammonia is concerned, we need mainly three elements:
\begin{enumerate}
    \item \emph{Motion}: Vibration of the N nucleus that goes up and down the plane of H nuclei;
    \item \emph{Associated potential}: The harmonic potential with a barrier centred on $z=0$. The N nucleus can penetrate the barrier via quantum tunnelling;
    \item \emph{States}: The two quantum states at lowest energy associated with this potential, indicated as + and -, are the equivalent of those of the barrierless harmonic potential, and each rotational state defined by quantum numbers (J, K) (NH3 is symmetric-top) is splitted in the two inversion states + and -
\end{enumerate}
\begin{figure}[h!]
    \centering 
    \includegraphics[width=0.6\linewidth]{img/inversiontransition.png}
    \caption{Development of energy level splitting when the potential barrier is reduced. (a) With barrier present, the wavefunctions have similar forms and energy levels are closely spaced; (b) as barrier is reduced, wavefunctions take on distinctly different form and energy levels separate.}
    \label{fgr:inversiontransition}
\end{figure}
The two inversion levels + and - are further split into hyperfine levels due to the non-zero spin of the N nucleus and of the three H nuclei. So, for example, in the rotational state $(j,k)=(1,1)$ the two inversion levels split into 5 hyperfine sub-levels; in total, the inversion transition (1,1) is splitted into 18 hyperfine components.

\subsubsection{Vibrational Levels}

In principle, we should also consider the energy level produced by vibration motions of biatomic molecules. Classically, we know that at large separations the Coulombian attraction of electrons on nuclei is dominating, while, on the shorter scales, prevails the Coulombian repulsion of nuclei. For biatomic molecules, the best approximation for this particular potential is the \emph{Morse potential}
\begin{equation}
    V(r) = D_{e}(1-\exp(-a(r-r_{e})))^{2}
    \label{eq:morsepotential}
\end{equation}
which we can Taylor-expand in a neighborhood of $r_{e}$, finding that the first non-vanishing order is the second one
\begin{equation}
    V(r) = \frac{1}{2}k(r-r_{e})^{2} +o(r^{3}) \quad k = \left(\parparder{V}{r}\right)_{r=r_{e}} = 2a^{2}D_{e}
\end{equation}
and so the Morse potential is equivalent to an harmonic potential, which we know fairly well how to treat, also quantistically speaking. Note that unharmonic contributions grow (significantly) relevant moving to larger distances from the equilibrium position.